\section{Effective closed sets and finite avoidance certificates}
\label{sec:effective}

The existence proof uses finite but extremely large random tables.
It also uses real parameters and an open-neighborhood repair.
Nevertheless, the resulting compact sets can be made effective by
searching directly for finite certificates. This route avoids having
to implement the entire probabilistic construction.

\subsection{A finite rational normalized certificate}

A \emph{rational periodic open set} is a finite union of rational open
arcs on $\R/\Z$, lifted periodically to $\R$. Its density is a
computable rational number, obtained by merging the arcs on one period.
An arc crossing zero is represented by an interval $(a,b)$ with
$0\le a<1$ and $a<b\le a+1$, together with its integer translates.
Dyadic rational endpoints suffice.

\begin{proposition}[Terminating search for normalized blockers]\label{prop:effective-H}
Given rational $0<\alpha<\beta<1$ and rational $0<p<1$, there is an algorithm
that returns a rational periodic open set $H$ and an integer $N\ge1$
such that $\rho(H)<6p$ and
\begin{equation}\label{eq:finite-H}
 \forall x\in[0,1]\ \forall t\in[1,2]\ \forall q\in[\alpha,\beta]
 \quad \bigvee_{n=1}^{N}(x+tq^n\in H).
\end{equation}
The analogous statement holds for a rational semialgebraic compact
parameter set and a computably given sequence of rational polynomials
satisfying \cref{thm:poly-family}. These hypotheses are promises;
the algorithm need not decide them from an arbitrary infinite sequence.
\end{proposition}

\begin{proof}
Enumerate pairs $(H,N)$ consisting of finite dyadic open arc unions
and positive integers, in an order eventually reaching every pair.
Discard candidates with density at least $6p$. For each remaining
candidate, decide \eqref{eq:finite-H} by quantifier elimination in
the ordered real field. The expressions $x+tq^n$ are polynomials with
rational coefficients. They lie in $[0,3]$, so membership in $H$
is a finite disjunction over only finitely many integer translates
of its arcs. The entire formula is therefore a first-order sentence
about rational polynomial inequalities. Real quantifier elimination
is effective; see, for example, \cite{PerrucciRoy}.

To prove termination, apply \cref{prop:normalized} with $p/2$,
obtaining an open periodic $H_0$ of density at most $3p$.
For each point of the compact parameter box there is a witnessing
integer $n$ and a small open interval, with dyadic endpoints, containing
the witnessing value and contained in $H_0$. By continuity of the
original expression, the same interval and exponent work on a
neighborhood of that parameter point. Finitely many of these
neighborhoods cover the compact box. Their finitely many witnessing
intervals project to a rational periodic open subset $H$ of $H_0$;
let $N$ be the largest witnessing exponent. Then
$\rho(H)\le3p<6p$ and \eqref{eq:finite-H} holds. Hence the search
must encounter a valid pair.

For polynomial families, replace $q^n$ by $f_n(z)$ and include the
rational semialgebraic formula defining $\Theta$ in the quantified
domain. Continuity, compactness, and effective real polynomial
arithmetic give the same proof.
\end{proof}

This is a proof of a total algorithm on the stated inputs, not a
reported computation of the full hitting set. The certificate search
may be enormous. A more direct implementation could enumerate the
finite tables, eliminate the residual-center quantifiers, and isolate
algebraic endpoints, but no such implementation is required for the
effective existence result.

\subsection{A computable measure modulus}

Use \cref{prop:effective-H} for every normalized family in the global
assembly. Index the families by $\ell\ge1$: this can mean the intervals
$I_\ell$ for the geometric theorem, the pairs $(d,j)$ for mixtures,
or the tuples used in \cref{cor:exp-polynomial,thm:signed-exp-poly}. In each case there
is a computable enumeration of rational semialgebraic families with
computably given rational polynomials. Use
\[
 p_{\ell,k}=\frac\eps{24}2^{-\ell}2^{-(k+1)}
 \quad(\ell\ge1,\ k\ge0).
\]
For $N\ge1$, let $C_N$ contain precisely the signed contractions with
$1\le\ell\le N$ and $0\le k<N$, and set
\begin{equation}\label{eq:KN}
 K_N=[0,1]\setminus C_N.
\end{equation}
The set $K_N$ is a computable finite union of closed rational intervals,
and the sequence is nested. The omitted density budget is bounded by
\begin{align}
 m(K_N)-m(K)
 &\le12\sum_{\ell>N\ \mathrm{or}\ k\ge N}p_{\ell,k}\notag\\
 &=\frac\eps2\left[1-(1-2^{-N})^2\right]
 \le\eps\,2^{-N}.
 \label{eq:measure-modulus}
\end{align}
Here the sums are over the union of the two omitted index ranges;
the overlap is not counted twice. This proves
\eqref{eq:effective-rate-intro}. Since $m(K_N)$ is rational and its
error has the computable bound \eqref{eq:measure-modulus}, $m(K)$
is a computable real. The same enumeration of rational intervals
shows that $K$ is effectively closed relative to $[0,1]$, and hence
as a compact subset of $\R$ after including its exterior in the
open complement.

A different rational accuracy can be computed independently. No
uniform operation on an unspecified noncomputable real $\eps$ is
asserted. For a nonrational positive accuracy, the existence statements
follow by choosing a smaller rational accuracy.

\subsection{Finite prefixes with a positive distance margin}

\begin{theorem}[Effective robust certificates]\label{thm:certificates}
Let $F$ be the effectively constructed periodic set from
\cref{thm:main}, and let $C=\R\setminus F$. Given rational parameters
\[
 R_0>0,\quad 0<\sigma<S_0,\quad 0<\alpha<\beta<1,
\]
and an integer $J\ge1$, there is an algorithm returning integers
$L,N\ge1$, $N\ge J$, and a rational $\delta>0$, together with a
finite rational interval description of $C_L$ modulo one, such that
\begin{equation}\label{eq:robust-certificate}
 \forall |x|\le R_0\ \forall s\in[-S_0,-\sigma]\cup[\sigma,S_0]
 \ \forall q\in[\alpha,\beta]
 \quad \max_{J\le n\le N}\dist(x+sq^n,F)\ge\delta.
\end{equation}
More precisely, one of these terms has its open $\delta$-neighborhood
contained in $C_L\subset C$.
The same assertion holds for $F^*_\eps$ and a fixed normalized
mixture family \eqref{eq:mixture-family}, with $q^n$ replaced by
$f_n(w,q)$.
\end{theorem}

\begin{proof}
Every orbit in the compact parameter domain has a hit in $C$ at an
index $n\ge J$, by the infinitely-many-hits property. That hit
belongs to some finite stage $C_L$. As $C_L$ is open, the point has
a positive distance from its complement. Continuity of the original
polynomial in the orbit parameters preserves a smaller positive
margin on a neighborhood of those parameters. Compactness gives
finitely many such neighborhoods. Take the largest stage and index
needed and a positive rational margin smaller than all their margins.
This proves existence of a finite certificate.

For effectivity enumerate $(L,N,\delta)$ with $N\ge J$ and positive
rational $\delta\le1$. On the bounded range of possible orbit values,
compute the finite list of rational open interval components of $C_L$,
using enough integer translates to include their $1$-neighborhoods.
A sufficient certificate is the sentence asserting that for every
parameter tuple some $J\le n\le N$ and some component $(a,b)$ satisfy
\[
 a+\delta<x+sq^n<b-\delta.
\]
It is decidable by real quantifier elimination. The existence argument
shows that a candidate eventually passes. The shrunken-interval
condition proves \eqref{eq:robust-certificate}. For normalized
mixtures the domain is compact semialgebraic and the terms are again
rational polynomials in the parameters, so the same reasoning applies.
\end{proof}

The strict shrunken-interval certificate is slightly stronger than a
bare distance inequality and easier to verify. It is finite,
polynomial, and independent of the infinite tail of the construction.
It therefore supplies a natural interface for eventual formalization.

\subsection{An interface suitable for ProveIt}

A formal development can separate three layers. The first is finite
combinatorics and probability: windows, key nesting, selectors, and the
conditional failure identity. The second is the finite geometric
certificate that all parameters in a box have a hit. The third is
Lebesgue measure and countable assembly, with the explicit tail bound
\eqref{eq:measure-modulus}. A proof assistant need not trust a floating-point
plot, an informal random sample, or a claim that an infinite list was
searched. Each finite certificate is a quantified polynomial statement
with exact rational data.

The present package supplies ordinary mathematical proofs and exact
Python checks. It does not include a Lean proof, an implementation of
real quantifier elimination, or an explicitly materialized full avoiding
set. These are distinct deliverables from the effective existence and
termination theorems proved above.
